\documentclass[a4paper,11pt]{article}
\usepackage{exam-style-842-unit}
\examinehead{842 信息化概论 · 第09单元练习册}{练习三\quad 农业信息化与数字乡村}
\begin{document}

\unitcover{842 信息化概论 · 单元练习册}{第09单元\quad 信息化与农业应用}{练习三\quad 农业信息化现状、数字乡村与综合论述（共 100 分）}{约 70 分钟}{我国农业信息化发展现状、主要问题与解决办法；数字乡村的内涵与实现途径；农业信息服务的内容与模式；物联网、大数据、人工智能、区块链在农业中的综合应用；本单元论述题的答题框架。}

\begin{knowmap}
\textbf{农业信息化内涵：}信息基础设施 · 信息资源 · 信息技术应用 · 信息服务 · 人才与政策环境\\[2pt]
\textbf{我国现状：}网络与平台基础明显改善、示范应用加快；但“最后一公里”、数据孤岛、人才与资金短板仍然突出\\[2pt]
\textbf{主要问题五维度：}基础设施 · 信息资源 · 人才 · 标准 · 资金（并表现为技术与生产脱节、区域发展不平衡）\\[2pt]
\textbf{数字乡村：}数字基础设施 + 智慧农业 + 乡村数字治理 + 信息惠民 + 网络文化 + 数字生态\\[2pt]
\textbf{农业信息服务：}政策法规·生产技术·市场行情·气象预警·金融保险·培训咨询；政府主导型·企业市场型·合作组织型·平台型\\[2pt]
\textbf{新技术综合应用：}物联网（感知与控制）· 大数据（分析与预测）· 人工智能（识别与决策）· 区块链（溯源与信任）\\[2pt]
\textbf{论述框架：}破题 $\rightarrow$ 背景 $\rightarrow$ 技术 $\rightarrow$ 应用 $\rightarrow$ 问题 $\rightarrow$ 对策 $\rightarrow$ 升华
\end{knowmap}

\begin{keyterms}
\textbf{农业信息化}：在农业全过程中广泛应用信息技术、开发利用农业信息资源、发展农业信息服务，提高农业生产、经营、管理与服务水平的过程。\\
\textbf{数字乡村}：伴随网络化、信息化和数字化在农业农村经济社会发展中的应用，以及农民现代信息技能提高而内生的农业农村现代化发展和转型进程。\\
\textbf{《数字乡村发展战略纲要》}：2019 年中共中央办公厅、国务院办公厅印发，是数字乡村建设的纲领性文件。\\
\textbf{农业信息服务}：向农民、合作社和涉农企业提供政策法规、生产技术、市场行情、气象预警、金融保险等信息与咨询服务的活动。\\
\textbf{益农信息社}：信息进村入户工程在村级设立的服务站点，提供公益服务、便民服务、电子商务和培训体验服务。\\
\textbf{12316}：全国农业系统公益服务统一专用号码，为农民提供政策、技术、市场等咨询服务。\\
\textbf{信息进村入户}：以村级信息服务站为载体，把信息服务延伸到村、到户的农业信息化重点工程。\\
\textbf{农业农村大数据}：农业农村领域多源海量数据的汇聚、治理、分析与应用，是农业智能决策的数据基础。\\
\textbf{数字鸿沟}：不同地区、群体在信息基础设施接入和信息技术应用能力上的差距。\\
\textbf{区块链}：一种去中心化、不可篡改、可追溯的分布式账本技术，可在多方之间建立信任。\\
\textbf{农产品质量安全溯源}：借助物联网、区块链等记录农产品生产、加工、流通全过程信息，实现来源可查、去向可追、责任可究。\\
\textbf{农业电子商务}：利用电子商务平台开展农产品与农业生产资料交易、农业服务对接的经营方式。\\
\textbf{乡村数字治理}：运用数字技术开展数字党建、网格化管理、政务服务、应急管理等，提升乡村治理效能。\\
\textbf{乡村振兴战略}：按照产业兴旺、生态宜居、乡风文明、治理有效、生活富裕的总要求推进农业农村现代化的战略。
\end{keyterms}

\begin{contrast}{农业信息化 / 数字乡村（每行依次为：核心内涵·覆盖范围·建设目标）}
\citem{农业信息化}{信息技术在农业全过程的推广应用}{主要覆盖农业生产、经营、管理与服务}{提高农业生产效率与管理服务水平}
\citem{数字乡村}{数字技术赋能乡村生产、生活、治理、生态}{覆盖农业、农村、农民及乡村治理与公共服务}{以数字化推动乡村全面振兴}
\end{contrast}

\begin{contrast}{农业信息服务的三种模式（每行依次为：供给主体·主要形式·特点与局限）}
\citem{政府主导型}{农技推广机构、农业主管部门}{农技站咨询、12316 热线、信息进村入户}{公益性强、覆盖面广，但内容时效性与针对性有限}
\citem{企业市场型}{电信运营商、电商与平台企业}{短信彩信、农业 APP、电商平台}{机制灵活、服务多样，但收费与盈利模式制约可持续性}
\citem{合作组织型}{合作社、行业协会、社会化服务组织}{统一采购、技术指导、产销对接}{贴近农户、针对性强，但规模与专业能力有限}
\end{contrast}

\begin{contrast}{物联网 / 区块链 / 人工智能在农业中的作用（每行依次为：技术特征·农业应用·解决的核心问题）}
\citem{物联网}{全面感知、可靠传输、智能控制}{环境监测、精准灌溉、设施控制、养殖监测}{解决“看不见、控不准”的问题}
\citem{区块链}{去中心化、不可篡改、可追溯}{农产品溯源、供应链金融、土地确权}{解决“信不过、追不到”的问题}
\citem{人工智能}{学习、识别、预测与决策}{病虫害图像识别、产量预测、智能农机}{解决“算不准、靠经验”的问题}
\end{contrast}

\qtypename{一、单项选择题}{10 题，每题 3 分，共 30 分}
\q 农业信息化是指（\quad）
\choices{A. 仅在农业机械上安装定位设备}{B. 信息技术在农业全过程和农村经济社会各领域的推广应用}
{C. 只建设农村宽带网络}{D. 只发展农产品电子商务}

\q 我国农业信息化中常说的“最后一公里”问题是指（\quad）
\choices{A. 农产品运输距离远}{B. 信息与技术服务难以到达农户和田间地头}
{C. 农村道路硬化不足}{D. 网络设备价格高}

\q 2019 年印发、作为数字乡村建设纲领性文件的是（\quad）
\choices{A.《数字乡村发展战略纲要》}{B.《国家信息化发展战略》}{C.《乡村振兴促进法》}{D.《“互联网+”行动计划》}

\q 数字乡村的内涵是（\quad）
\choices{A. 在农村普及互联网}{B. 以数字技术赋能农业农村现代化和乡村全面振兴}
{C. 只建设智慧农业园区}{D. 只推进农村电子商务}

\q 下列属于农业信息服务内容的是（\quad）
\choices{A. 政策法规、生产技术、市场行情与气象预警}{B. 仅农产品价格信息}
{C. 仅天气预报信息}{D. 仅农机购置补贴信息}

\q 信息进村入户工程在村级设立的服务站点称为（\quad）
\choicesline{A. 益农信息社}{B. 农技推广站}{C. 供销社}{D. 农村信用社}

\q 全国农业系统公益服务统一专用号码是（\quad）
\choicesline{A. 12315}{B. 12316}{C. 12345}{D. 12306}

\q 区块链在农业中最典型的应用是（\quad）
\choices{A. 农产品质量安全溯源}{B. 农田灌溉控制}{C. 农机作业调度}{D. 农业气象预报}

\q 农业物联网一般由感知层、网络层、支撑层和（\quad）组成。
\choicesline{A. 应用层}{B. 表示层}{C. 会话层}{D. 传输层}

\q 解决我国农业信息化问题的根本途径是（\quad）
\choices{A. 单纯增加硬件投入}{B. 政府主导、多方参与，基础设施、资源、人才、标准、服务协同推进}
{C. 完全依靠市场自发调节}{D. 照搬国外发展模式}

\qtypename{二、多项选择题}{3 题，每题 4 分，共 12 分；每题有两个或两个以上正确项}
\q 我国农业信息化存在的主要问题有（\quad）
\choices{A. 农村信息基础设施仍较薄弱}{B. 农业信息资源分散、共享不足}
{C. 农业信息化人才缺乏}{D. 农业机械化水平过高}

\q 数字乡村建设的主要内容包括（\quad）
\choices{A. 数字基础设施建设}{B. 智慧农业}{C. 乡村数字治理}{D. 信息惠民服务}

\q 区块链在农业中的应用包括（\quad）
\choices{A. 农产品质量安全溯源}{B. 农业供应链金融}{C. 农村土地确权登记}{D. 农田自动灌溉控制}

\qtypename{三、判断题}{6 题，每题 2 分，共 12 分；正确的画“√”，错误的画“×”并写出改正后的表述}
\q 农业信息化仅指在农业生产环节应用信息技术。\quad（\quad）\quad 改正：\fixblank

\q 数字乡村就是实现农村网络全覆盖。\quad（\quad）\quad 改正：\fixblank

\q 区块链具有去中心化、不可篡改、可追溯等特点，适合用于农产品质量安全溯源。\quad（\quad）\quad 改正：\fixblank

\q 益农信息社是信息进村入户工程在村级设立的服务站点。\quad（\quad）\quad 改正：\fixblank

\q 农业信息服务只向农民提供生产技术信息。\quad（\quad）\quad 改正：\fixblank

\q 数字乡村是乡村振兴的战略方向，是数字中国的重要内容。\quad（\quad）\quad 改正：\fixblank

\qtypename{四、填空题}{6 空，每空 1 分，共 6 分}
\q 2019 年印发的《\fillblanks[2.6cm]发展战略纲要》是数字乡村建设的纲领性文件。

\q 全国农业系统公益服务统一专用号码是\fillblanks[1.6cm]。

\q 信息进村入户工程在村级设立的服务站点称为\fillblanks[2.2cm]。

\q 农业物联网一般由感知层、网络层、支撑层和\fillblanks[1.8cm]组成。

\q 区块链具有去中心化、不可篡改、\fillblanks[2.0cm]等特点，可用于农产品溯源。

\q 农业信息服务的内容包括政策法规、生产技术、市场行情、\fillblanks[2.6cm]和金融保险等。

\newpage
\qtypename{五、简答题}{3 题，每题 8 分，共 24 分}
\q 简述我国农业信息化发展的现状及存在的主要问题。\anslines{2.4cm}

\q 针对我国农业信息化存在的主要问题，提出相应的解决办法（对策）。\anslines{2.4cm}

\q 什么是数字乡村？简述数字乡村的内涵与实现途径。\anslines{2.4cm}

\newpage
\qtypename{六、论述题}{1 题，共 16 分}
\q 结合乡村振兴战略和数字乡村建设，论述物联网、大数据、人工智能、区块链等信息技术在我国农业生产与乡村治理中的应用，分析当前农业信息化存在的主要问题并提出对策。要求：观点明确、层次清楚，须结合水稻、柑橘、生猪或茶叶等具体产业举例，并体现“技术—应用—问题—对策”的逻辑。\anslines{6cm}

\end{document}
